%=== Câu 1 ===%
\begin{ex}
% ID: L12C1B2-34-TN
% Mức: NB
	{\color{red}\footnotesize\ttfamily [L12C1B1-TN-24]}\par %HienID
	Một chất lỏng có khối lượng $m$, nhiệt hoá hơi riêng là $L$, nhiệt nóng chảy riêng là $\lambda$, nhiệt dung riêng $c$. Nhiệt hóa hơi có thể được xác định thông qua công thức nào sau đây?
	\choice
	{$Q = \Delta U - A$}
	{\True $Q=L.m$}
	{$Q=mc\Delta t$}
	{$Q=\lambda m$}
	\loigiai{
		Nhiệt hóa hơi được xác định bằng biểu thức $Q=Lm$.
	}
\end{ex}

%=== Câu 2 ===%
\begin{ex}
% ID: L12C1B2-35-TN
% Mức: NB
	{\color{red}\footnotesize\ttfamily [L12C1B1-TN-28]}\par %HienID
	Nhiệt lượng cần truyền cho một kilôgam chất lỏng để nó hóa hơi hoàn toàn ở nhiệt độ xác định gọi là 
	\choice
	{nhiệt nóng chảy riêng}
	{nhiệt hóa hơi}
	{\True nhiệt hóa hơi riêng}
	{nhiệt dung riêng} 
	\loigiai{
		Nhiệt lượng cần truyền cho một kilôgam chất lỏng để nó hóa hơi hoàn toàn ở nhiệt độ xác định cho biết nhiệt
		hóa hơi riêng của chất lỏng đó.
	}
\end{ex}

%=== Câu 3 ===%
\begin{ex}
% ID: L12C1B2-36-TN
% Mức: NB
	{\color{red}\footnotesize\ttfamily [L12C1B1-TN-30]}\par %HienID
	Đơn vị của nhiệt hoá hơi riêng là
	\choice
	{$kg$}
	{\True $J/kg$}
	{$J$}
	{$kg/J$}
	\loigiai{
		Đơn vị của nhiệt hoá hơi riêng là ${J} / {kg}$.
		
	}
\end{ex}

%=== Câu 4 ===%
\begin{ex}
% ID: L12C1B2-37-TN
% Mức: VD
	{\color{red}\footnotesize\ttfamily [L12C1B1-TN-31]}\par %HienID
	Biết nhiệt độ sôi, nhiệt dung riêng và nhiệt hóa hơi riêng của nước lần lượt là $100^\circ \mathrm{C}$, $4\,200\ \mathrm{J/kg}.K$ và $2{,}3.10^6\ \mathrm{J/kg}$. Nhiệt lượng cần cung cấp để làm hóa hơi hoàn toàn $2\ \mathrm{kg}$ nước từ nhiệt độ $20^\circ \mathrm{C}$ là
	\choice
	{$2{,}6.10^6\ \mathrm{J}$}
	{\True $5{,}3.10^6\ \mathrm{J}$}
	{$26{,}10^6\ \mathrm{J}$}
	{ $53{,}10^6\ \mathrm{J}$}
	\loigiai{
		Nhiệt lượng cần cung cấp để làm hóa hơi hoàn toàn 2 kg nước từ nhiệt độ $20^{\circ} {C}$ là:\\
		${Q}={mc} \Delta {t}+{mL}=2 . 4\,200 .(100-20)+2 . 2{,}3 . 10^{6}=5\,272\,000 {~J} \approx 5{,}3 . 10^{6} {~J}$.
		
	}
\end{ex}

%=== Câu 5 ===%
\begin{ex}
% ID: L12C1B2-38-TN
% Mức: VD
	{\color{red}\footnotesize\ttfamily [L12C1B1-TN-32]}\par %HienID
	
	\begin{wrapfigure}{r}{4cm}
		\vspace{-10pt}
		\centering
		% TODO: \usepackage{graphicx} required
		% TODO: \usepackage{graphicx} required
		\includegraphics[width=3.7cm,keepaspectratio]{"ImagesGPT/hơi nước nóng"}
		\vspace{-10pt}
	\end{wrapfigure}
	
	Biết nhiệt hóa hơi riêng của nước là $L = 2{,}3 \cdot 10^6\ \mathrm{J/kg}$. Nhiệt lượng cần cung cấp để làm bay hơi hoàn toàn $100\ \mathrm{g}$ nước ở nhiệt độ sôi là
	
	\choice
	{$23 \cdot 10^5\ \mathrm{J}$}
	{\True $23 \cdot 10^4\ \mathrm{J}$}
	{$23 \cdot 10^7\ \mathrm{J}$}
	{$23 \cdot 10^6\ \mathrm{J}$}
	\loigiai{
		
		Đổi $m = 100\ \mathrm{g} = 0{,}1\ \mathrm{kg}$. Áp dụng công thức tính nhiệt lượng hóa hơi: $Q = m \cdot L$, ta có:
		\[Q = 0{,}1 \cdot 2{,}3 \cdot 10^6 = 23 \cdot 10^4\ \mathrm{J}.\]
		
	}
\end{ex}

%=== Câu 6 ===%
\begin{ex}
% ID: L12C1B2-39-DS
% Mức: VD
	Một đứa bé bị sốt ở $38{,}2\ °\mathrm{C}$ được cho dùng thuốc hạ sốt. Sau khi uống thuốc, cơn sốt hạ xuống $36{,}7\ °\mathrm{C}$. Giả sử bay hơi mồ hôi là cách duy nhất để giảm nhiệt độ cơ thể. Dữ kiện: khối lượng $m = 25\ \mathrm{kg}$; nhiệt dung riêng $c = 1\,000\ \mathrm{cal/(kg \cdot °C)}$; nhiệt hoá hơi của nước $L = 580\ \mathrm{cal/g}$.
	\choiceTF
	{\True Độ giảm nhiệt độ là $1{,}5\ \mathrm{K}$.}
	{Nhiệt lượng tỏa ra là $37\,500\ \mathrm{J}$.}
	{Khối lượng nước bay hơi là $48\ \mathrm{g}$.}
	{\True Nếu $13\%$ nhiệt tỏa ra mất vào môi trường, khối lượng nước bay hơi là khoảng $56{,}25\ \mathrm{g}$.}
	\loigiai{
		\begin{itemchoice}
			\item $\Delta T = 38{,}2 - 36{,}7 = 1{,}5\ °\mathrm{C} = 1{,}5\ \mathrm{K}$. Mệnh đề đúng.
			
			\item $Q = m \cdot c \cdot \Delta T = 25 \cdot 1\,000 \cdot 1{,}5 = 37\,500\ \mathrm{cal}$ (không phải $\mathrm{J}$). Mệnh đề sai.
			
			\item Từ $Q = m_{\mathrm{nước}} \cdot L$: $m_{\mathrm{nước}} = \frac{Q}{L} = \frac{37\,500}{580} \approx 64{,}7\ \mathrm{g} \neq 48\ \mathrm{g}$. Mệnh đề sai.
			
			\item Với $87\%$ nhiệt được dùng: $m_{\mathrm{nước}} = \frac{0{,}87 \cdot 37\,500}{580} = 56{,}25\ \mathrm{g}$. Mệnh đề đúng.
		\end{itemchoice}
	}
\end{ex}

%=== Câu 7 ===%
\begin{bt}
% ID: L12C1B2-40-TLN
% Mức: VD
	{\color{red}\footnotesize\ttfamily [L12C1B1-TLN-7]}\par %HienID
	Tính nhiệt lượng cần thiết để làm nóng chảy hoàn toàn $2\ \mathrm{kg}$ nước đá ở $0^\circ\mathrm{C}$. Biết nhiệt nóng chảy riêng của nước đá là $3{,}4 \cdot 10^5\ \mathrm{J/kg}$. Tính kết quả theo đơn vị $\mathrm{kJ}$.
	\shortans[oly]{$680$}
	\loigiai{
		\[
		\begin{aligned}
			Q &= m\lambda \\
			&= 2 \cdot 3{,}4 \cdot 10^5 \\
			&= 6{,}8 \cdot 10^5\ \mathrm{J} = 680\ \mathrm{kJ}.
		\end{aligned}
		\]
	}
\end{bt}

%=== Câu 8 ===%
\begin{bt}
% ID: L12C1B2-41-TLN
% Mức: VD
	{\color{red}\footnotesize\ttfamily [L12C1B1-TLN-8]}\par %HienID
	Một khối chất có khối lượng $0{,}5\ \mathrm{kg}$ nhận nhiệt lượng $150\ \mathrm{kJ}$ để hóa hơi hoàn toàn ở nhiệt độ sôi. Tính nhiệt hóa hơi riêng của chất này theo đơn vị $10^5\ \mathrm{J/kg}$ (làm tròn đến một chữ số thập phân).
	\shortans[oly]{$3{,}0$}
	\loigiai{
		\[
		\begin{aligned}
			L &= \frac{Q}{m} \\
			&= \frac{150\,000}{0{,}5} \\
			&= 300\,000\ \mathrm{J/kg} = 3{,}0 \cdot 10^5\ \mathrm{J/kg}.
		\end{aligned}
		\]
	}
\end{bt}

%=== Câu 9 ===%
\begin{bt}
% ID: L12C1B2-42-TLN
% Mức: VD
	{\color{red}\footnotesize\ttfamily [L12C1B1-TLN-9]}\par %HienID
	Để làm bay hơi hoàn toàn $50\ \mathrm{g}$ nước ở $100^\circ\mathrm{C}$, cần cung cấp một nhiệt lượng bằng bao nhiêu, biết nhiệt hóa hơi riêng của nước là $2{,}3 \cdot 10^6\ \mathrm{J/kg}$? Tính kết quả theo đơn vị $\mathrm{kJ}$.
	\shortans[oly]{$115$}
	\loigiai{
		\[
		\begin{aligned}
			Q &= mL \\
			&= 0{,}05 \cdot 2{,}3 \cdot 10^6 \\
			&= 115\,000\ \mathrm{J} = 115\ \mathrm{kJ}.
		\end{aligned}
		\]
	}
\end{bt}

%=== Câu 10 ===%
\begin{bt}
% ID: L12C1B2-43-TL
% Mức: VDC
	{\color{red}\footnotesize\ttfamily [L12C1B1-TL-4]}\par %HienID
	Một khối nước đá có khối lượng $300\ \mathrm{g}$ ở $0^\circ\mathrm{C}$ được đun nóng để nóng chảy hoàn toàn thành nước ở $0^\circ\mathrm{C}$, sau đó tiếp tục đun để nước nóng lên đến $100^\circ\mathrm{C}$ và cuối cùng hóa hơi hoàn toàn ở $100^\circ\mathrm{C}$. Biết nhiệt nóng chảy riêng của nước đá là $3{,}34 \cdot 10^5\ \mathrm{J/kg}$, nhiệt dung riêng của nước là $4200\ \mathrm{J/(kg\cdot K)}$, nhiệt hóa hơi riêng của nước là $2{,}3 \cdot 10^6\ \mathrm{J/kg}$.
	\begin{enumerate}
		\item Giải thích tại sao trong giai đoạn nóng chảy và giai đoạn hóa hơi, nhiệt độ của nước không đổi dù vẫn được cung cấp nhiệt.
		\item Tính nhiệt lượng cần thiết cho từng giai đoạn: nóng chảy, tăng nhiệt độ từ $0^\circ\mathrm{C}$ lên $100^\circ\mathrm{C}$, và hóa hơi hoàn toàn.
		\item Tính tổng nhiệt lượng cần cung cấp cho toàn bộ quá trình.
	\end{enumerate}
	\loigiai{
		Đổi đơn vị: $m = 300\ \mathrm{g} = 0{,}3\ \mathrm{kg}$.
		\begin{enumerate}
			\item Trong giai đoạn nóng chảy và hóa hơi, toàn bộ nhiệt lượng cung cấp được dùng để phá vỡ (hoặc làm suy yếu) lực liên kết giữa các phân tử nhằm chuyển từ thể này sang thể khác (tăng khoảng cách và thế năng tương tác giữa các phân tử), chứ không dùng để làm tăng động năng chuyển động nhiệt của các phân tử, nên nhiệt độ giữ nguyên không đổi trong các giai đoạn này.
			\item Nhiệt lượng cần cho giai đoạn nóng chảy:
			\[
			\begin{aligned}
				Q_1 &= m\lambda \\
				&= 0{,}3 \cdot 3{,}34 \cdot 10^5 \\
				&= 100\,200\ \mathrm{J} \approx 100{,}2\ \mathrm{kJ}.
			\end{aligned}
			\]
			Nhiệt lượng cần để tăng nhiệt độ từ $0^\circ\mathrm{C}$ lên $100^\circ\mathrm{C}$:
			\[
			\begin{aligned}
				Q_2 &= mc\Delta t \\
				&= 0{,}3 \cdot 4200 \cdot 100 \\
				&= 126\,000\ \mathrm{J} = 126\ \mathrm{kJ}.
			\end{aligned}
			\]
			Nhiệt lượng cần cho giai đoạn hóa hơi hoàn toàn:
			\[
			\begin{aligned}
				Q_3 &= mL \\
				&= 0{,}3 \cdot 2{,}3 \cdot 10^6 \\
				&= 690\,000\ \mathrm{J} = 690\ \mathrm{kJ}.
			\end{aligned}
			\]
			\item Tổng nhiệt lượng cần cung cấp:
			\[
			\begin{aligned}
				Q &= Q_1 + Q_2 + Q_3 \\
				&= 100{,}2 + 126 + 690 \\
				&= 916{,}2\ \mathrm{kJ}.
			\end{aligned}
			\]
		\end{enumerate}
	}
\end{bt}

%=== Câu 11 ===%
\begin{bt}
% ID: L12C1B2-44-TL
% Mức: VDC
	{\color{red}\footnotesize\ttfamily [L12C1B1-TL-5]}\par %HienID
	Người ta đun nóng liên tục một khối nước đá khối lượng $1\ \mathrm{kg}$ ban đầu ở $-10^\circ\mathrm{C}$ bằng một nguồn nhiệt có công suất không đổi $\mathcal{P} = 2\ \mathrm{kW}$, cho đến khi toàn bộ nước đá tan chảy hết thành nước ở $0^\circ\mathrm{C}$. Biết nhiệt dung riêng của nước đá là $2100\ \mathrm{J/(kg\cdot K)}$, nhiệt nóng chảy riêng của nước đá là $3{,}34 \cdot 10^5\ \mathrm{J/kg}$. Bỏ qua hao phí nhiệt ra môi trường.
	\begin{enumerate}
		\item Mô tả đồ thị nhiệt độ theo thời gian của quá trình trên: chỉ rõ đoạn nào ứng với giai đoạn nước đá tăng nhiệt độ và đoạn nào ứng với giai đoạn nóng chảy.
		\item Tính thời gian cần thiết để nước đá tăng nhiệt độ từ $-10^\circ\mathrm{C}$ lên $0^\circ\mathrm{C}$.
		\item Tính thời gian cần thiết để nước đá nóng chảy hoàn toàn ở $0^\circ\mathrm{C}$.
		\item Tính tổng thời gian của toàn bộ quá trình trên.
	\end{enumerate}
	\loigiai{
		\begin{enumerate}
			\item Đồ thị nhiệt độ - thời gian của quá trình gồm hai đoạn: đoạn thứ nhất là một đường dốc lên từ $-10^\circ\mathrm{C}$ đến $0^\circ\mathrm{C}$, ứng với giai đoạn nước đá tăng nhiệt độ (còn ở thể rắn); đoạn thứ hai là một đường nằm ngang tại $0^\circ\mathrm{C}$, ứng với giai đoạn nước đá đang nóng chảy (nhiệt độ không đổi).
			\item Nhiệt lượng cần để nước đá tăng nhiệt độ từ $-10^\circ\mathrm{C}$ lên $0^\circ\mathrm{C}$:
			\[
			\begin{aligned}
				Q_1 &= mc\Delta t \\
				&= 1 \cdot 2100 \cdot 10 \\
				&= 21\,000\ \mathrm{J}.
			\end{aligned}
			\]
			Thời gian tương ứng:
			\[
			\begin{aligned}
				t_1 &= \frac{Q_1}{\mathcal{P}} \\
				&= \frac{21\,000}{2\,000} \\
				&= 10{,}5\ \mathrm{s}.
			\end{aligned}
			\]
			\item Nhiệt lượng cần để nước đá nóng chảy hoàn toàn:
			\[
			\begin{aligned}
				Q_2 &= m\lambda \\
				&= 1 \cdot 3{,}34 \cdot 10^5 \\
				&= 334\,000\ \mathrm{J}.
			\end{aligned}
			\]
			Thời gian tương ứng:
			\[
			\begin{aligned}
				t_2 &= \frac{Q_2}{\mathcal{P}} \\
				&= \frac{334\,000}{2\,000} \\
				&= 167\ \mathrm{s}.
			\end{aligned}
			\]
			\item Tổng thời gian của toàn bộ quá trình:
			\[
			\begin{aligned}
				t &= t_1 + t_2 \\
				&= 10{,}5 + 167 \\
				&= 177{,}5\ \mathrm{s}.
			\end{aligned}
			\]
		\end{enumerate}
	}
\end{bt}

%=== Câu 12 ===%
\begin{ex}
% ID: L12C1B2-45-DS
% Mức: NB
Cho các phát biểu sau về sự thay đổi nội năng và định luật 1 của nhiệt động lực học, hãy xác định đúng hay sai.
\choiceTF
{\True (P1) Khi vật nhận nhiệt lượng $Q$ và thực hiện công $A$, độ biến thiên nội năng của vật là $\Delta U = Q - A$}
{(P2) Nội năng của vật tăng khi vật nhận công và tỏa nhiệt ra môi trường với độ lớn nhiệt lượng tỏa ra lớn hơn công mà vật nhận}
{\True (P3) Truyền nhiệt và thực hiện công là hai cách duy nhất để làm thay đổi nội năng của một hệ}
{(P4) Theo quy ước, khi vật nhận công thì giá trị $A$ trong công thức $\Delta U = A + Q$ là âm}
\loigiai{
\begin{itemchoice}
\item (Đúng) (P1) Khi vật nhận nhiệt lượng $Q$ và thực hiện công $A$, độ biến thiên nội năng của vật là $\Delta U = Q - A$. (Vì): $
\item (Sai) (P2) Nội năng của vật tăng khi vật nhận công và tỏa nhiệt ra môi trường với độ lớn nhiệt lượng tỏa ra lớn hơn công mà vật nhận.
\item (Đúng) (P3) Truyền nhiệt và thực hiện công là hai cách duy nhất để làm thay đổi nội năng của một hệ.
\item (Sai) (P4) Theo quy ước, khi vật nhận công thì giá trị $A$ trong công thức $\Delta U = A + Q$ là âm.
\end{itemchoice}
}
\end{ex}

%=== Câu 13 ===%
\begin{ex}
% ID: L12C1B2-46-TN
% Mức: NB
Thả một miếng sắt đã được nung nóng vào một cốc nước lạnh thì
\choice
{nội năng của cốc nước và miếng sắt đều giảm}
{\True nội năng của cốc nước tăng, nội năng của miếng sắt giảm}
{nội năng của cốc nước giảm, nội năng của miếng sắt tăng}
{nội năng của cốc nước và miếng sắt đều tăng}
\loigiai{
Khi thả miếng sắt nóng vào cốc nước lạnh, nhiệt độ của sắt giảm trong khi nhiệt độ của nước tăng.
 Điều này dẫn đến nội năng của miếng sắt giảm và nội năng của nước tăng lên cho đến khi đạt cân bằng nhiệt.
 Quá trình này là minh chứng cho sự trao đổi nhiệt giữa hai vật có nhiệt độ chênh lệch.
 Phân tích quá trình truyền nhiệt: nội năng truyền từ vật có nhiệt độ cao sang vật có nhiệt độ thấp, đây là sự cân bằng nội năng thông qua truyền nhiệt.
}
\end{ex}

%=== Câu 14 ===%
\begin{ex}
% ID: L12C1B2-47-DS
% Mức: NB
Chỉ ra câu đúng, sai trong các câu sau:
\choiceTF
{\True (P1) Động năng của các phân tử trong một khối khí xác định là như nhau}
{(P2) Thế năng của mỗi phân tử khí trong một bình kín là giống nhau}
{(P3) Nội năng của một khối khí không liên quan tới tổng năng lượng của các nguyên tử tạo thành khối khí đó}
{(P4) Nội năng của một khối khí phụ thuộc vào nhiệt độ và thể tích của khối khí đó}
\loigiai{
\begin{itemchoice}
\item (Đúng) (P1) Động năng của các phân tử trong một khối khí xác định là như nhau. (Vì): (P1) Đúng. Động năng trung bình của các phân tử trong một khối khí xác định là như nhau ở một nhiệt độ nhất định
\item (Sai) (P2) Thế năng của mỗi phân tử khí trong một bình kín là giống nhau. (Vì): (P2) Sai. Thế năng tương tác giữa các phân tử khí là rất nhỏ và không cố định, không phải là giống nhau cho mỗi phân tử
\item (Sai) (P3) Nội năng của một khối khí không liên quan tới tổng năng lượng của các nguyên tử tạo thành khối khí đó. (Vì): (P3) Sai. Nội năng của một khối khí là tổng động năng và thế năng của các phân tử cấu tạo nên khối khí đó, do đó nó liên quan đến tổng năng lượng của các nguyên tử
\item (Sai) (P4) Nội năng của một khối khí phụ thuộc vào nhiệt độ và thể tích của khối khí đó. (Vì): (P4) Sai. Đối với khí lí tưởng, nội năng chỉ phụ thuộc vào nhiệt độ. Đối với khí thực, nội năng phụ thuộc vào cả nhiệt độ và thể tích. Phát biểu này không hoàn toàn đúng trong mọi trường hợp (đối với khí lí tưởng
\end{itemchoice}
}
\end{ex}

%=== Câu 15 ===%
\begin{ex}
% ID: L12C1B2-48-TN
% Mức: NB
Quá trình làm thay đổi nội năng của vật bằng cách cho nó tiếp xúc với vật khác khi
\choice
{nhiệt độ của chúng bằng nhau gọi là sự trao đổi công}
{có sự chênh lệch nhiệt độ giữa chúng gọi là sự nhận công}
{\True có sự chênh lệch nhiệt độ giữa chúng gọi là sự truyền nhiệt}
{nhiệt độ của chúng bằng nhau gọi là sự truyền nhiệt}
\loigiai{
Quá trình truyền nhiệt diễn ra khi có sự chênh lệch nhiệt độ giữa hai vật. Năng lượng sẽ chuyển từ vật có nhiệt độ cao sang vật có nhiệt độ thấp cho đến khi cân bằng nhiệt
}
\end{ex}

%=== Câu 16 ===%
\begin{ex}
% ID: L12C1B2-49-TN
% Mức: NB
Phương pháp nào sau đây \textbf{không} làm tăng nội năng của vật?
\choice
{Nước trong nồi được đun nóng}
{Cọ xát miếng kim loại vào mặt bàn}
{Viên bi được thả vào nước nóng}
{\True Viên bi rơi trong chân không}
\loigiai{
Trong chân không, không có không khí nên không có lực cản hay ma sát. Viên bi rơi tự do mà không có sự trao đổi nhiệt lượng hay thực hiện công lên viên bi, do đó nội năng của viên bi không thay đổi.
Thả rơi viên bi trong môi trường chân không, nội năng không thay đổi
}
\end{ex}

%=== Câu 17 ===%
\begin{ex}
% ID: L12C1B2-50-TN
% Mức: NB
Nội năng của một vật là
\choice
{nhiệt lượng vật nhận được trong quá trình truyền nhiệt}
{tổng nhiệt lượng và cơ năng mà vật nhận được trong quá trình truyền nhiệt và thực hiện công}
{tổng động năng và thế năng của vật}
{\True tổng động năng và thế năng của các phân tử cấu tạo nên vật}
\loigiai{
Tổng động năng và thế năng của các phân tử cấu tạo nên vật
}
\end{ex}

%=== Câu 18 ===%
\begin{ex}
% ID: L12C1B2-51-TN
% Mức: NB
Trong trường hợp nào dưới đây, sự biến đổi nội năng của vật \textbf{không} phải do thực hiện công?
\choice
{Mài dao}
{Đóng đinh}
{Khuấy nước}
{\True Nung sắt trong lò}
\loigiai{
Nung sắt trong lò là quá trình truyền nhiệt, nhiệt năng từ lò truyền sang sắt làm tăng nhiệt độ và nội năng của sắt mà không cần thực hiện công.
 Quá trình thực hiện công làm cho nội năng của vật thay đổi. Vật nhận công thì nội năng tăng, vật thực hiện công cho vật khác thì nội năng giảm.
Nung sắt trong lò là quá trình truyền nhiệt, nhiệt năng từ lò truyền sang sắt làm tăng nhiệt độ và nội năng của sắt mà không cần thực hiện công
}
\end{ex}

%=== Câu 19 ===%
\begin{ex}
% ID: L12C1B2-52-TN
% Mức: NB
Để xảy ra sự truyền nhiệt khi cho hai vật tiếp xúc với nhau thì hai vật phải có
\choice
{thể tích khác nhau}
{chiều cao khác nhau}
{\True nhiệt độ khác nhau}
{khối lượng khác nhau}
\loigiai{
Khi cho hai vật chênh lệch nhiệt độ tiếp xúc nhau, năng lượng nhiệt luôn truyền từ vật có nhiệt độ cao hơn sang vật có nhiệt độ thấp hơn. Quá trình truyền nhiệt kết thúc khi hai vật ở cùng nhiệt độ (trạng thái cân bằng nhiệt)
}
\end{ex}

%=== Câu 20 ===%
\begin{ex}
% ID: L12C1B2-53-TN
% Mức: NB
Một người dùng tay cầm một viên nước đá đang tan, năng lượng nhiệt truyền
\choice
{\True từ tay người cầm sang viên nước đá}
{từ tay người cầm sang viên nước đá và ngược lại}
{từ viên nước đá sang tay người cầm}
{không xảy ra quá trình truyền nhiệt}
\loigiai{
Năng lượng nhiệt truyền từ tay người cầm sang viên nước đá do sự chênh lệch nhiệt độ.
Quá trình truyền nhiệt diễn ra từ nơi có nhiệt độ cao hơn (tay) sang nơi có nhiệt độ thấp hơn (viên nước đá)
}
\end{ex}

%=== Câu 21 ===%
\begin{ex}
% ID: L12C1B2-54-TN
% Mức: NB
Hai vật rắn $A$ và $B$ được đặt ở gần nhau. Vật $A$ có nhiệt độ cao hơn vật $B$. Nội dung nào sau đây là \textbf{không} chính xác?
\choice
{Tốc độ trung bình của các phân tử trong vật $B$ nhỏ hơn tốc độ trung bình của các phân tử trong vật $A$}
{Năng lượng nhiệt được truyền từ vật $A$ sang vật $B$}
{Quá trình truyền nhiệt giữa hai vật dừng lại khi hai vật có nhiệt độ bằng nhau}
{\True Nếu nhiệt độ của hai vật nhỏ hơn nhiệt độ của môi trường xung quanh thì không có sự truyền nhiệt giữa hai vật}
\loigiai{
Truyền nhiệt xảy ra khi có chênh lệch nhiệt độ, bất kể nhiệt độ môi trường cao hay thấp. Vì vậy, nếu hai vật có nhiệt độ chênh lệch, quá trình trao đổi nhiệt vẫn diễn ra
}
\end{ex}

%=== Câu 22 ===%
\begin{ex}
% ID: L12C1B2-55-TN
% Mức: NB
Theo định luật $I$ của nhiệt động lực học, hệ thức $\Delta U = A + Q$ khi $Q \gt  0$ và $A \lt  0$ mô tả quá trình
\choice
{hệ nhận nhiệt và nhận công}
{hệ truyền nhiệt và sinh công}
{hệ truyền nhiệt và sinh công}
{\True hệ nhận nhiệt và sinh công}
\loigiai{
Theo định luật $I$ của nhiệt động lực học, công thức:
 $\Delta U = A + Q$
 Khi $Q \gt  0$, hệ nhận nhiệt; khi $A \lt  0$, hệ thực hiện công ra ngoài.
Trong quá trình này, nội năng của hệ tăng lên, tương ứng với hệ nhận nhiệt và sinh công
}
\end{ex}

%=== Câu 23 ===%
\begin{ex}
% ID: L12C1B2-56-TN
% Mức: NB
Một vật đang được làm nóng sao cho thể tích của vật không thay đổi. Nội năng của vật
\choice
{không thay đổi}
{giảm đi}
{\True tăng lên}
{tăng lên rồi giảm đi}
\loigiai{
Khi làm nóng một vật mà thể tích không thay đổi, toàn bộ nhiệt lượng nhận được sẽ làm tăng nội năng. Vì không có sự giãn nở, tất cả năng lượng được tích tụ vào nội năng
}
\end{ex}

%=== Câu 24 ===%
\begin{ex}
% ID: L12C1B2-57-TN
% Mức: NB
Phát biểu nào sau đây về nội năng là \textbf{không} đúng?
\choice
{\True Nội năng là nhiệt lượng vật nhận được trong quá trình truyền nhiệt}
{Nội năng của một vật có thể tăng lên hoặc giảm đi}
{Nội năng có thể chuyển hóa thành các dạng năng lượng khác}
{Nội năng của vật bao gồm tổng động năng và thế năng của các phân tử cấu tạo nên vật}
\loigiai{
Phần năng lượng nhiệt trao đổi giữa hai vật gọi là nhiệt lượng. Trong quá trình truyền nhiệt, không có sự chuyển hóa năng lượng từ dạng này sang dạng khác mà chỉ có sự chuyển nội năng từ vật này sang vật khác.
Có sự chênh lệch nhiệt độ giữa chúng gọi là sự truyền nhiệt. Khi hai vật có nhiệt độ khác nhau tiếp xúc với nhau thì xảy ra quá trình truyền nhiệt. Quá trình này làm thay đổi nội năng của vật.
phần năng lượng nhiệt trao đổi giữa hai vật gọi là nhiệt lượng .Trong quá trình truyền nhiệt không có sự chuyển hoá năng lượng từ dạng này sang dạng khác mà chỉ có sự truyển nội năng từ vật này sang vật khác có sự chênh lệch nhiệt độ giữa chúng gọi là sự truyền nhiệt
}
\end{ex}

%=== Câu 25 ===%
\begin{ex}
% ID: L12C1B2-58-TN
% Mức: NB
Một vật được làm nóng sao cho thể tích của vật không thay đổi thì nội năng của vật
\choice
{giảm}
{\True tăng}
{giảm rồi tăng}
{không thay đổi}
\loigiai{
Theo định luật nhiệt động lực học, nội năng của một vật phụ thuộc vào nhiệt độ và thể tích.
Trong trường hợp này, vật được **làm nóng** nên nhiệt độ của vật **tăng lên**.
Do thể tích không thay đổi, không có sự giãn nở, nên toàn bộ nhiệt lượng cung cấp sẽ làm tăng nội năng.
Áp dụng công thức:
 $\Delta U = nC_V \Delta T$
 Vì $\Delta T \gt  0$, nên $\Delta U \gt  0$, tức là nội năng của vật tăng lên
}
\end{ex}

%=== Câu 26 ===%
\begin{ex}
% ID: L12C1B2-59-TN
% Mức: NB
Nội năng của một vật
\choice
{chỉ phụ thuộc vào nhiệt độ của vật}
{chỉ phụ thuộc thể tích của vật}
{\True phụ thuộc thể tích và nhiệt độ của vật}
{không phụ thuộc thể tích và nhiệt độ của vật}
\loigiai{
Nội năng của một vật phụ thuộc thể tích và nhiệt độ của vật. Khi nhiệt độ hoặc thể tích thay đổi, nội năng cũng thay đổi.
Nội năng của vật phụ thuộc vào nhiệt độ $T$ và thể tích $V$ của vật. Khi truyền nhiệt hoặc thực hiện công, nội năng sẽ thay đổi.
Quá trình thực hiện công hoặc truyền nhiệt đều có thể làm thay đổi nội năng của vật. Khi hai vật có nhiệt độ khác nhau tiếp xúc nhau, quá trình truyền nhiệt sẽ xảy ra
}
\end{ex}

%=== Câu 27 ===%
\begin{ex}
% ID: L12C1B2-60-TN
% Mức: NB
Hệ thức định luật $I$ nhiệt động lực học $\Delta U=A+Q$. Khi $Q<0$ và $A>0$ mô tả quá trình
\choice
{vật dẫn nhiệt và sinh công}
{vật truyền nhiệt và sinh công}
{vật dẫn nhiệt và nhận công}
{\True vật truyền nhiệt và nhận công}
\loigiai{
Hệ thức định luật $I$ nhiệt động lực học khi $Q<0$ và $A>0$ mô tả quá trình vật truyền nhiệt và nhận công.
Ghi rõ thêm nội dung về dấu như:
 
 
A. $Q>0$: Hệ nhận nhiệt
 
B. $Q<0$: Hệ truyền nhiệt
 
C. $A>0$: Hệ sinh công
 
D. $A<0$: Hệ nhận công
}
\end{ex}

%=== Câu 28 ===%
\begin{ex}
% ID: L12C1B2-61-TN
% Mức: NB
Hai vật có nhiệt độ khác nhau tiếp xúc với nhau. Nhiệt năng được truyền từ
\choice
{vật có nội năng lớn sang vật có nội năng nhỏ hơn}
{\True vật có nhiệt độ cao sang vật có nhiệt độ thấp hơn}
{vật có kích thước lớn sang vật có kích thước nhỏ}
{vật có khối lượng lớn sang vật có khối lượng nhỏ hơn}
\loigiai{
Khi hai vật có nhiệt độ khác nhau tiếp xúc, năng lượng nhiệt luôn truyền từ vật có nhiệt độ cao hơn sang vật có nhiệt độ thấp hơn, quá trình này kết thúc khi đạt trạng thái cân bằng nhiệt
}
\end{ex}

%=== Câu 29 ===%
\begin{ex}
% ID: L12C1B2-62-TN
% Mức: NB
Trường hợp nào sau đây thể hiện nội năng của một vật bị thay đổi thông qua quá trình truyền nhiệt?
\choice
{Mũi tiện tạo ra nhiệt khi tiện phôi}
{Dầu que diêm nóng lên khi đánh diêm}
{Xoa tay khiến tay ta nóng lên}
{\True Nước xà phòng làm nguội mũi tiện}
\loigiai{
Nội năng của một vật có thể thay đổi thông qua hai cách:
 
 
• **Truyền nhiệt**: nhiệt năng chuyển từ vật nóng sang vật lạnh.
 
• **Thực hiện công**: ma sát, nén ép hoặc va chạm.
 
 Trong trường hợp này, nước xà phòng làm nguội mũi tiện là một quá trình truyền nhiệt, không phải là thực hiện công
}
\end{ex}

%=== Câu 30 ===%
\begin{ex}
% ID: L12C1B2-63-TN
% Mức: NB
Nội năng của một vật phụ thuộc vào
\choice
{khối lượng và thể tích của vật}
{khối lượng của vật}
{khối lượng và nhiệt độ của vật}
{\True nhiệt độ và thể tích của vật}
\loigiai{
Nội năng của một vật là hàm phụ thuộc vào nhiệt độ và thể tích của vật, ký hiệu là:
 $U = f(T, V)$
 Khi nhiệt độ hoặc thể tích thay đổi, nội năng của vật cũng sẽ thay đổi theo
}
\end{ex}

%=== Câu 31 ===%
\begin{ex}
% ID: L12C1B2-64-TN
% Mức: NB
Một vật trượt từ đỉnh xuống chân mặt phẳng nghiêng. Trong quá trình này, nội năng của vật
\choice
{\True tăng}
{giảm}
{không thay đổi}
{giảm rồi tăng}
\loigiai{
Khi một vật trượt từ đỉnh xuống chân mặt phẳng nghiêng, lực ma sát giữa vật và mặt phẳng nghiêng làm cho một phần cơ năng (động năng và thế năng) chuyển hóa thành nội năng.
Điều này khiến cho nhiệt độ của vật và mặt phẳng nghiêng tăng lên, đồng nghĩa với việc nội năng của vật cũng tăng lên.
Quá trình này thể hiện rõ sự chuyển hóa năng lượng từ cơ năng sang nội năng thông qua ma sát.
Nội năng của một vật là tổng nội năng và thế năng tương tác của các phân tử cấu tạo nên vật. Khi có ma sát, một phần cơ năng chuyển hóa thành nội năng làm nhiệt độ vật tăng lên
}
\end{ex}

%=== Câu 32 ===%
\begin{ex}
% ID: L12C1B2-65-TN
% Mức: NB
Một viên đạn bắn vào một khối gỗ đặt trên một mặt phẳng nằm ngang nhẵn thì toàn bộ động năng của viên đạn chuyển hóa thành
\choice
{động năng của hệ (đạn + gỗ)}
{nội năng của khối gỗ}
{\True nội năng của hệ (đạn + gỗ)}
{động năng và nội năng của hệ (đạn + gỗ)}
\loigiai{
Khi viên đạn bắn vào khối gỗ, do mặt phẳng nằm ngang là nhẵn (bỏ qua ma sát), toàn bộ động năng của viên đạn sẽ chuyển hóa thành nội năng của hệ (đạn + gỗ).
Điều này xảy ra do sự ma sát, va chạm và biến dạng trong quá trình đạn đi vào khối gỗ, làm tăng nhiệt độ và nội năng của cả hệ.
Không có sự chuyển hóa thành động năng sau va chạm vì hệ dừng lại hoặc di chuyển rất ít (không có ma sát để bảo toàn động năng).
Phụ thuộc thể tích và nhiệt độ của vật. Khi hai vật có nhiệt độ khác nhau tiếp xúc với nhau thì xảy ra quá trình truyền nhiệt. Quá trình này làm thay đổi nội năng của các vật
}
\end{ex}

%=== Câu 33 ===%
\begin{ex}
% ID: L12C1B2-66-TN
% Mức: NB
Gọi $A$ và $Q$ lần lượt là công và nhiệt lượng mà hệ nhận được. Độ biến thiên nội năng của hệ được tính bằng công thức nào sau đây?
\choice
{\True $\Delta U = A + Q$}
{$\Delta U = A - Q$}
{$\Delta U = Q - A$}
{$\Delta U = -A - Q$}
\loigiai{
Độ biến thiên nội năng của hệ được tính theo công thức:
 $\Delta U = A + Q$
 Đây là công thức tổng quát của định luật $I$ nhiệt động lực học
}
\end{ex}

%=== Câu 34 ===%
\begin{ex}
% ID: L12C1B2-67-TN
% Mức: NB
Quy ước về dấu nào sau đây phù hợp với định luật $I$ của nhiệt động lực học $\Delta U=Q+A$ khi vật truyền nhiệt và nhận công?
\choice
{\True $A> 0$ và $Q< 0$}
{$A< 0$ và $Q> 0$}
{$A> 0$ và $Q> 0$}
{$A< 0$ và $Q< 0$}
\loigiai{
Khi vật truyền nhiệt thì $Q< 0$, và vật nhận công nên $A> 0$.
Ý nghĩa về dấu các đại lượng:
 
 
A. $Q>0$: Vật nhận nhiệt lượng.
 
B. $Q<0$: Vật truyền nhiệt lượng.
 
C. $A>0$: Vật nhận công.
 
D. $A<0$: Vật thực hiện công.
}
\end{ex}

%=== Câu 35 ===%
\begin{ex}
% ID: L12C1B2-68-TN
% Mức: NB
Bỏ một đồng xu được nung nóng vào một cốc nước đá đang tan, trước khi trạng thái cân bằng nhiệt được thiết lập thì nội năng của đồng xu
\choice
{tăng lên rồi giảm đi}
{\True giảm đi}
{tăng lên}
{không thay đổi}
\loigiai{
Khi bỏ một đồng xu đã được nung nóng vào cốc nước đá, quá trình truyền nhiệt xảy ra từ đồng xu (nhiệt độ cao) sang nước đá (nhiệt độ thấp).
Quá trình này làm cho nhiệt độ của đồng xu giảm xuống, dẫn đến nội năng của nó cũng giảm theo.
Đây là quá trình truyền nhiệt theo nguyên tắc cân bằng nhiệt: nhiệt độ sẽ san bằng cho đến khi đạt trạng thái cân bằng
}
\end{ex}

%=== Câu 36 ===%
\begin{ex}
% ID: L12C1B2-69-TN
% Mức: NB
Phát biểu nào sau đây là đúng khi nói về nội năng?
\choice
{Các vật khác nhau ở cùng nhiệt độ thì có cùng nội năng}
{Khi các vật tăng tốc độ thì thế năng của các phân tử tăng và nội năng cũng tăng}
{\True Thực hiện công và truyền nhiệt thì có thể làm thay đổi nội năng của vật}
{Động năng và thế năng hấp dẫn của vật là một phần nội năng của vật}
\loigiai{
Nội năng của một vật có thể thay đổi thông qua hai quá trình chính:
 
 
• Thực hiện công: lực tác động lên vật làm biến đổi nội năng.
 
• Truyền nhiệt: trao đổi nhiệt lượng giữa hai vật có chênh lệch nhiệt độ.
 
 Cả hai quá trình này đều dẫn đến sự thay đổi nội năng của vật
}
\end{ex}

%=== Câu 37 ===%
\begin{ex}
% ID: L12C1B2-70-TN
% Mức: NB
$\Delta U$ là độ biến thiên nội năng của một vật, $Q$ là nhiệt lượng vật trao đổi với môi trường, vật thực hiện hoặc nhận một công $A$. Công thức tổng quát của định luật $I$ nhiệt động lực học là
\choice
{$\Delta U = A$}
{$\Delta U = Q$}
{$Q + A = 0$}
{\True $\Delta U = A + Q$}
\loigiai{
Công thức tổng quát của nguyên lý một nhiệt động lực học:
 $\Delta U = A + Q$
 Độ biến thiên nội năng của một vật bằng tổng công và nhiệt lượng mà vật nhận được
}
\end{ex}

%=== Câu 38 ===%
\begin{ex}
% ID: L12C1B2-71-TN
% Mức: NB
Khi nội năng của một vật tăng lên thì năng lượng trung bình của các phân tử cấu tạo nên vật cũng tăng lên. Khi đó
\choice
{có động năng của các phân tử tăng lên}
{chỉ có thế năng của các phân tử tăng lên}
{động năng của các phân tử chắc chắn tăng lên còn thế năng của chúng có thể thay đổi không đáng kể}
{\True động năng và thế năng của các phân tử chắc chắn cùng tăng lên}
\loigiai{
- Nội năng của vật là tổng động năng và thế năng tương tác của các phân tử. Khi nội năng tăng, cả động năng và thế năng của các phân tử đều có thể tăng. Cụ thể hơn, động năng trung bình của phân tử liên quan đến nhiệt độ, và thế năng tương tác giữa các phân tử liên quan đến khoảng cách giữa chúng (thể tích). Nếu nội năng tăng, thường là do nhiệt độ tăng (tức động năng tăng) hoặc do thay đổi thể tích (tức thế năng tăng), hoặc cả hai.
 Kết quả: $\boxed{D}$
}
\end{ex}

%=== Câu 39 ===%
\begin{ex}
% ID: L12C1B2-72-TN
% Mức: NB
Hãy tìm câu sai trong các câu sau: Để làm thay đổi nội năng của một vật, ta
\choice
{cung cấp nhiệt lượng cho vật}
{\True thực hiện công nhất vật lên theo phương vuông góc với mặt đất một đoạn $1\,\text{m}$}
{cho vật trượt từ độ cao $1\,\text{m}$ xuống mặt đất bằng mặt phẳng nghiêng với góc nghiêng 60° so với mặt đất}
{cho vật truyền nhiệt lượng sang một vật khác có nhiệt độ thấp hơn}
\loigiai{
- Có hai cách làm thay đổi nội năng của vật: thực hiện công hoặc truyền nhiệt.
 - Phương án $A$ (cung cấp nhiệt lượng) và $D$ (truyền nhiệt lượng) là các hình thức truyền nhiệt.
 - Phương án $C$ (cho vật trượt từ độ cao xuống) là hình thức thực hiện công (công của lực ma sát làm tăng nội năng của vật và mặt phẳng).
 - Phương án $B$ (thực hiện công nhất vật lên) làm thay đổi thế năng cơ học của vật trong trường trọng lực, chứ không làm thay đổi trực tiếp nội năng của vật (trong điều kiện lí tưởng, không có ma sát). Do đó, $B$ là câu sai.
 Kết quả: $\boxed{B}$
}
\end{ex}

%=== Câu 40 ===%
\begin{ex}
% ID: L12C1B2-73-TN
% Mức: NB
Phát biểu nào sau đây là \textbf{đúng} về nội năng của một vật?
\choice
{Nội năng của một vật chỉ phụ thuộc vào nhiệt độ của vật đó}
{Nội năng của một vật không thể thay đổi nếu không có sự truyền nhiệt}
{\True Nội năng là tổng động năng của các phân tử cấu tạo nên vật và thế năng tương tác giữa chúng}
{Đơn vị của nội năng là Calo}
\loigiai{
A. (A) **Sai** vì nội năng của một vật phụ thuộc vào cả nhiệt độ và thể tích của vật. Đối với khí lí tưởng, nội năng mới chỉ phụ thuộc vào nhiệt độ.
 
B. (B) **Sai** vì nội năng có thể thay đổi thông qua thực hiện công hoặc truyền nhiệt.
 
C. (C) **Đúng** vì nội năng của một vật là tổng động năng và thế năng tương tác của các phân tử cấu tạo nên vật.
 
D. (D) **Sai** vì đơn vị của nội năng trong hệ SI là Joule (J).
}
\end{ex}

